\section{Related work}\label{sec:related}

\subsection{Classical quantum simulation}

Qiskit Aer~\cite{Qiskit} provides exact state-vector evolution up to
approximately 30 qubits on commodity CPUs and matrix-product-state
(MPS) compression up to approximately 50 qubits at constrained bond
dimension. Cirq~\cite{Cirq} is the Google equivalent; both share a
unitary-evolution semantics and differ primarily in gate-set defaults
and measurement collapse routines. \qmirror{} uses Aer as the default
engine with a Cirq alternative scaffolded for Google-native gate sets.

The pragmatic gap that \qmirror{} closes is the \textbf{measurement
collapse sampling}: Aer and Cirq default to a pseudo-random Mersenne
Twister or similar for CDF inversion. \qmirror{} substitutes ANU QRNG
bytes (or HMAC-DRBG keyed by ANU) at the CDF-inversion step. This is
the only substantive substrate-honesty change; everything upstream of
measurement is already exact.

\subsection{NIST statistical RNG validation}

NIST SP~800-22 Rev.~1a~\cite{NIST-SP-800-22} specifies a 15-test
statistical battery for cryptographic RNG validation. Tier-1 tests
(monobit, frequency-within-block, runs, longest-run-of-ones,
DFT-spectral, non-overlapping-template, overlapping-template,
Maurer-universal, linear-complexity, serial, approximate-entropy,
cumulative-sums-forward, cumulative-sums-backward, random-excursions,
random-excursions-variant) are applied to bit strings of length
$n = 10^{6}$ per test with 10 sequences per test. $p$-values must
exceed $0.01$.

\qmirror{} cond.4 runs tier-1+ (the seven tier-1 tests plus serial and
approximate-entropy from tier-2) on HMAC-DRBG bytes seeded by ANU
QRNG. PASS criterion: all 7+ tests at $p > 0.01$.

\subsection{Integrated Information Theory 4.0}

Tononi et al.\ IIT~4.0~\cite{Tononi-IIT4} defines \phistar{}
(intrinsic information) as the minimum-information partition (MIP)
over the substrate's transition probability matrix (TPM). The
\texttt{pyphi} reference implementation~\cite{pyphi} computes
\phistar{} via the \texttt{sia()} function on the
\texttt{feature/iit-4.0} branch (commit \texttt{b78d0e3} lineage). For
systems with $n \le 6$ nodes \texttt{pyphi} runs full MIP search; for
$6 < n \le 12$ it uses the \texttt{CUT\_ONE\_APPROXIMATION}
heuristic.

\qmirror{} cond.6 wraps \texttt{pyphi.sia()} via subprocess isolation
(\texttt{\_python\_bridge/iit\_mip\_runner.py}) and validates
byte-identical \phistar{} output against four reference TPMs from the
\texttt{braket\_iit40\_mip\_2026\_05\_02} cycle
(\texttt{and\_ionq\_forte1}, \texttt{maj\_ionq\_forte1},
\texttt{and\_sv1}, \texttt{maj\_sv1}; all four
\texttt{HONEST\_NEGATIVE} $\phistar = 0.0$).

\subsection{CHSH Bell inequality}

Clauser-Horne-Shimony-Holt~\cite{CHSH-1969} specifies the
four-correlator witness
$S = E(a,b) - E(a,b') + E(a',b) + E(a',b')$
with classical bound $\lvert S\rvert \le 2$ and quantum (Tsirelson)
bound $\lvert S\rvert \le 2\sqrt{2} \approx 2.828$~\cite{Tsirelson-1980}.
The canonical Aspect~\cite{Aspect-1982} experimental geometry uses
$R_{y}(-\theta)$ rotations with $\theta \in \{0, \pi/2, \pi/4,
-\pi/4\}$. \qmirror{}'s \texttt{chsh.run} implements this canonical
geometry and reproduces $S = 2.838$ at $n_{\text{trials}} = 1000$
against the IonQ Aria-1 reference $S = 2.808$ from
\texttt{nexus\_chsh\_bell\_2026\_05\_02}
($\dS = 0.030$, within $\pm 0.05$ band).
